\chapter*{函数章节重点题目回顾} % {第3章校本作业}

\section*{概念类}
\subsection*{函数的概念}
\begin{enumerate}
    \item 设函数$f(x)=\frac{1}{\sqrt{ax^2+ax+1}}$分别满足下列条件, 求实数$a$的取值范围: 
\begin{enumerate}
\item 定义域是$\mathbf{R}$; 
\item 值域是$(0,+\infty)$. 
\end{enumerate}
\begin{jieda}
\begin{enumerate}
    \item $ax^2+ax+1>0$恒成立. \\
    $a=0$时显然成立; \\
    $a\neq 0$时, $a>0$且$ \Delta<0$, 即$a\in(0,4)$. \\
    故$a\in[0,4)$. 
    
    \item 问题等价于$y=ax^2+ax+1$的值域包含$(0,+\infty)$. \\
    $a=0$显然不符合条件; \\
    $a\neq 0$时, $a>0$且$\Delta\geq0$, 即$a\geq 4$. \\
    故$a\geq 4$. \\
    \ps 学生容易犯两类错误：\\
        1. 误认为当$\sqrt{ax^2+ax+1}$的定义域是$\mathbf{R}$时，其值域是$(0, +\infty)$，从而错误地将问题转化为$\sqrt{ax^2+ax+1}$的定义域是$\mathbf{R}$，要求对应的二次函数开口向上且$\Delta < 0$，得到$[0, 4)$\\
        2. 计算$\Delta \geq 0$之后忘记开口方向的约束，得到$(-\infty, 0] \cup [4, +\infty)$
    \end{enumerate}
    \end{jieda}
    \item 函数$lg(ax^2+ax+1)$的值域为$\mathbf{R}$，求参数$a$的取值范围
    \begin{jieda}
        问题等价于$y=ax^2+ax+1$的值域包含$(0,+\infty)$. \\
        $a=0$显然不符合条件; \\
        $a\neq 0$时, $a>0$且$\Delta\geq0$, 即$a\geq 4$. \\
        故$a\geq 4$. \\
        \ps 道理同上一题，定义域不是$\mathbf{R}$不影响值域是$\mathbf{R}$.
    \end{jieda}
    \item 已知$f(x)=x+2, x\in[1,4]$, 则$y=f(x^2)+f^2(x)$的最大值是\blkda{22}. 
    \begin{jieda}
    先看定义域: $f(x)$的定义域为$[1,4]$，因此$f(x^2)$中$x^2\in[1,4]$, 故$f(x^2)$的定义域为$x\in[-2,2]$. 而$f^2(x)$的定义域与$f(x)$的定义域相同.\\
    两者相加，定义域取交集，因此$y=f(x^2)+f^2(x)$的定义域为$[1, 2]$ \\
    于是$y=x^2+2+(x+2)^2=2x^2+4x+6=2(x+1)^2+4\in[12,22]$. \\
    \ps 如果抽象函数求定义域、函数的运算求定义域仍然迷糊，请及时联系你的老师.
    \end{jieda}
    \item $y = x + \frac{4}{x+3} (x \in (-3, 0]$的值域为\blkda{$[1, +\infty)$}.
    \begin{jieda}
        换元$t = x+3$，$y = t+\frac{4}{t}-3 (t \in (0, 3])$，根据对勾函数性质可以得到值域为$[1, +\infty)$ \\
        \ps 出错方式或者是不理解渐近线，想象不到函数的值域到$+\infty$，带端点得到$[1, \frac{4}{3}]$，或者是换元时没有减3，得到$[4, +\infty)$；对于前者，熟悉函数的图像是基本要求，分析清楚单调性才能带入端点值是“普世价值”，而对于后者，请再仔细些.
    \end{jieda}
    \item 函数$f(x)=\frac{3^x}{9^x+1}-\frac{1}{2}$的值域是\blkda{$\left(-\frac{1}{2},0\right]$}. 
    \begin{jieda}
    因为$y=3^x+3^{-x}\in[2,+\infty)$, \\
    故$f(x)=\frac{1}{3^x+3^{-x}}-\frac{1}{2}\in\left(-\frac{1}{2},0\right]$. 
    \end{jieda}
    \item 若函数$f(x)=x^2-3x-4$的定义域为$[0,m]$, 值域为$\left[-\frac{25}{4},-4\right]$, 则实数$m$的取值范围是\blkda{$[\frac{3}{2},3]$}. 
    \begin{jieda}
    $f(x)=x^2-3x-4=(x-\frac{3}{2})^2-\frac{25}{4}$，因此$x = \frac{3}{2}$是函数值$-\frac{25}{4}$对应的唯一自变量，必然要在定义域内. 又因为$f(0)=f(3)=-4$, 是$f(x)=-4$唯二的解. 于是由二次函数的图像易得$m\in[\frac{3}{2},3]$. 
    \end{jieda}
    \item 已知函数$y=4^x-3\cdot2^x+3$的值域为$\left[1,7\right]$, 则$x$的范围可以是\dotfill(\kh{D})
    \xx{$\left[2,4\right]$}
    {$(-\infty,0)$}
    {$(0,1)\cup\left[2,4\right]$}
    {$\left(-\infty,0\right]\cup\left[1,2\right]$}
    \begin{jieda}
    $t^2-3t+3=1\iff t=1$或2, $t^2-3t+3=7\iff t=-1$或4. \\
    由于$2^x\in(0,+\infty)$, 故结合$y=t^2-3t+3$的图像可知, 必须要定义域必须满足有$2^x\in[2,4]\iff x\in[1,2]$. \\
    至于$[-1,1]$, $2^x$不能遍历, 所以$x\to-\infty$方向是随意的. 
    \end{jieda}
    \item 求下列函数的值域: 
    \begin{enumerate}
    \item $y=\frac{3+x}{4-x}$; 
    \item $y=\frac{2x^2-2x+3}{x^2-x+1}$; 
    \item $y=\frac{x^4+x^2+5}{(x^2+1)^2}$; 
    \item $y=x+4\sqrt{1-x}$; 
    \item $y=\frac{\sqrt{x^2+1}}{x-1}$. 
    \item $\sqrt{1+x} - \sqrt{2+x}$.
    \end{enumerate}
    \begin{jieda}
    \begin{enumerate}
    \item $y=-1+\frac{7}{4-x}\in(-\infty,-1)\cup(-1,+\infty)$. 
    
    \item $y=2+\frac{1}{x^2-x+1}$. \\
    因为$x^2-x+1\in[\frac{3}{4},+\infty)$, 故$y\in(2,\frac{10}{3}]$. 
    
    \item $y=\frac{x^4+2x^2+1+4-x^2}{x^4+2x^2+1}=1+\frac{4-x^2}{x^4+2x^2+1}$. \\
    令$t=4-x^2$. 当$t\in(-\infty,0)\cup(0,4]$, 则$t+\frac{25}{t}-10\in(-\infty,-20]\cup[\frac{1}{4},+\infty)$. \\
    于是当$t\leq 4$时, $y=\frac{t}{(t-5)^2}\in[-\frac{1}{20},4]$. \\
    所以$y=1+\frac{4-x^2}{x^4+2x^2+1}\in[\frac{19}{20},5]$. \\
    \ps 可化为$ax+\frac{b}{x}$的. 
    
    \item 令$t=\sqrt{1-x}\in[0,+\infty)$, 则$x=1-t^2$, 故$y=1-t^2+4t=-(t-2)^2+5\in(-\infty,5]$. \\
    \ps 可化为二次的. 
    
    \item 当$x>1$时, $y=\sqrt{\frac{x^2+1}{x^2-2x+1}}=\sqrt{1+\frac{2x}{x^2-2x+1}}=\sqrt{1+\frac{2}{x+\frac{1}{x}-2}}\in(1,+\infty)$. \\
    当$x<1$时, $y=-\sqrt{\frac{x^2+1}{x^2-2x+1}}=-\sqrt{1+\frac{2x}{x^2-2x+1}}$. \\
    再分为$(-\infty,0)\cup\{0\}\cup(0,1)$讨论可得$y\in(-\infty,-\frac{1}{\sqrt{2}}]$. \\
    综上, $y\in(-\infty,-\frac{1}{\sqrt{2}}]\cup(1,+\infty)$. %\\
    % \ps 可化为$ax+\frac{b}{x}$的. \\
    % \ps 对于$f(x)=\frac{x}{x^2+1}$而言, 在$(-\infty,-1]$上$\searrow$, 在$[-1,1]$上$\nearrow$, 在$[1,+\infty)$上$\searrow$. \\
    \item $y = \sqrt{1+x} - \sqrt{2+x} = \frac{(\sqrt{1+x} + \sqrt{2+x})(\sqrt{1+x} - \sqrt{2+x})}{(\sqrt{1+x} + \sqrt{2+x})} = \frac{-1}{(\sqrt{1+x} + \sqrt{2+x})}$ \\
    易知函数在定义域$[-1, +\infty)$内单调递增，最小值在左端点$x = -1$取$y = -1$，当$x$趋于无穷时，函数值趋于0。因此值域为$[-1, 0)$
    \end{enumerate}
    \end{jieda}
    \item 下列两组函数中表示同一函数的有\blkda{\ding{172}}. \\
    \begin{enumerate*}
        \item $y = \frac{\sqrt{1-x^2}}{|x+2|}$ 和 $y = \frac{\sqrt{1-x^2}}{x+2}$
        \item $y = \sqrt{x-1} \cdot \sqrt{x-2}$和$y = \sqrt{x^2-3x+2}$
    \end{enumerate*}
    \begin{jieda}
        本质上还是在考察定义域.
    \end{jieda}
\end{enumerate}
\subsection*{数形结合}
\begin{enumerate}
    \item 已知函数$f(x)=\begin{cases}
    2^{-x}-1, & x < 0\\
    \sqrt{x}, & x\geq 0
    \end{cases}$, 
    则使得$f(x) > 1$的$x$的取值范围是\blkda[3.5]{$(-\infty,-1) \cup (1,+\infty)$}. 
    \begin{jieda}
    画出函数图像就可以明确解集的形式.
    \end{jieda}
    \item 若方程$\sqrt{2x+1} = x+a$有两个不同的实数解时，则实数$a$的取值范围是\blkda{$\left[\frac{1}{2}, 1\right)$}.
    \begin{jieda}
        方程有两个不同的实数解可以转化为函数$y = \sqrt{2x+1}$和函数$y = x+a$拥有两个不同的交点.通过绘制函数图像可知，随着$a$的增大，$y = x+a$会逐渐向上平移，在变化的过程中，两个边界的情况分别是$y = x+a$过$(-\frac{1}{2}, 0)$和$y = x+a$与$y = \sqrt{2x+1}$相切(前者符合题意，后者不符合). \\
        当$y = x+a$过$(-\frac{1}{2}, 0)$的时候$a = \frac{1}{2}$；对原方程两边平方后要求$\Delta = 0$等价于两者相切，计算可得到$a = 1$. \\
        因此$a$的取值范围是$\left[\frac{1}{2}, 1\right)$. \\
        \ps 从函数的角度看方程，从函数图像运动的角度理解参数变化的影响.
    \end{jieda}
    \item 如果点$(m, n) (m \neq n)$是函数$f(x) = \sqrt{ax+2}(a < 0)$与其反函数$f^{-1}(x)$图像上的公共点，求$a$的取值范围
    \begin{jieda}
        $f^{-1}(x) = \frac{x^2-2}{a} (x \geq 0)$，由点$(m, n) (m \neq n)$是公共点可知$f(m) = n, f^{-1}(m) = n$，根据反函数的性质可知$f(n) = m, f^{-1}(n) = m$ \\
        选取$f^{-1}(m) = n, f^{-1}(n) = m$构建方程组（方便计算），即$\left \{ \begin{array}{ll}
             \frac{m^2-2}{a} = n  \\
             \frac{n^2-2}{a} = m
        \end{array}\right. \quad (m \geq 0, n \geq 0)$
        两式相减有$\frac{m^2-n^2}{a} = n-m$，即$\frac{(m+n)(m-n)}{a} = n-m$，因此$m = n$或$m+n = -a$，$m = n$对应着原函数与反函数在$y = x$的交点，必然存在，而$m+n = -a$对应着题目所求的不在$y = x$上的交点，不一定存在（可以通过带回$\frac{m^2-2}{a} = n$检验判断) \\
        因为$(m, n), (n, m)$均在$f(x)$上，故可以看做是直线$y = -(x-m)+n$与$f(x)$的交点，代入$m+n = -a$即$y = -x-a$与$f(x)$有两个交点. 从图像上可以判断出两个边界分别是$y = -x-a$与$f(x)$相切或者$y = -x-a$过$f(x)$的端点$(-\frac{2}{a}, 0)$. \\
        前者对应着方程$ax+2 = x^2+2ax+a^2$判别式为0，即$a^2-4(a^2-2)= 0$，解得$a = -\sqrt{\frac{8}{3}}$；后者对应着$\frac{2}{a} = a$，解得$a = -\sqrt{2}$，因此$a \in (-\sqrt{\frac{8}{3}}, -\sqrt{2}]$
    \end{jieda}
    
    \item 函数$g(x)=-\frac{x}{x+1}$的图像可以由函数$f(x)=\frac{x}{x-1}$的图像经过下列变换得到\dotfill(\kh{B})
    \xx{向左平移2个单位, 向上平移2个单位;}
    {向左平移2个单位, 向下平移2个单位;}
    {向右平移2个单位, 向上平移2个单位;}
    {向右平移2个单位, 向下平移2个单位.}
    \begin{jieda}
    $g(x)=-1+\frac{1}{x+1}$, $f(x)=1+\frac{1}{x-1}$, 故$g(x)=f(x+2)-2$. 
    \end{jieda}
    \item 函数$y=2^{1-x}+2$的图像可以由函数$y=(\frac{1}{2})^x$的图像经过下面变换得到\dotfill(\kh{C})
    \xx{向左平移1个单位, 向上平移2个单位}
    {向左平移1个单位, 向下平移2个单位}
    {向右平移1个单位, 向上平移2个单位}
    {向右平移1个单位, 向下平移2个单位}
    \begin{jieda}
    令$f(x)=(\frac{1}{2})^x=2^{-x}$, 则$y=f(x-1)+2=2^{1-x}+2$. \\
    另外, 用$f(0)=1$也是可以的, 关键在于``若你明白了什么是代表元素那就自由的''. 
    \end{jieda}
    \item 函数$y=\abs{x^2-1}+1$的图像与函数$y=\lg(x+10)$的图像交点个数为\blkda{2}. 
    \item 已知函数$f(x)=2^x-2$, 则函数$y=\abs{f(\abs{x})}$的图像只可能是\dotfill(\kh{A})
    \begin{minipage}{\linewidth}
    \centering
    \includegraphics[totalheight=3cm]{./bodyxb/src/34713785651095890.png}
    \end{minipage}
    \item 函数$y = \frac{x^2}{2^x-2^{-x}}$的图象大致为\dotfill(\kh{A})
    \begin{figure}[h]
        \centering
        \includegraphics[width=0.6\linewidth]{./bodyxb/src/q35.png}
    \end{figure}
    \begin{jieda}
        首先可以判断出来该函数是奇函数，并且当$x > 0$时$y >0$。通过带入特值($x=1, x=\frac{1}{2}$)（或求导）等可以判断出来函数在$x > 0$时不是单调递减的，因此选A
    \end{jieda}
    % \item 函数$y = (|x|+1)ln|x|$的图象大致为\dotfill(\kh{B})
    % \begin{figure}[h]
    %     \centering
    %     \includegraphics[width=0.6\linewidth]{./bodyxb/src/q34.png}
    % \end{figure}
    % \begin{jieda}
    %     首先可以判断出来该函数为偶函数，其次可以判断出来在$x > 0$时函数是单调递增的，因此选B。
    % \end{jieda}
    \item 已知关于$x$的方程$3^{-x} = \frac{2a+3}{5-a}$有正的实根时，实数$a$的取值范围是\blkda[3]{$(-\frac{3}{2}, \frac{2}{3})$}.
    \begin{jieda}
        首先根据指数函数的图像将问题转化为$\frac{2a+3}{5-a} \in (0, 1)$，进一步分离常数后有$\frac{13}{5-a} \in (2, 3)$，做倒数处理有$\frac{5-a}{13} \in (\frac{1}{3}, \frac{1}{2})$，所以$5-a \in (\frac{13}{3}, \frac{13}{2})$，故$a \in (-\frac{3}{2}, \frac{2}{3})$
    \end{jieda}
    \item 关于$x$的方程$|3^x-1| = k$只有一个解，则$k$的取值范围是\blkda[3]{$[1, +\infty) \cup \{0\}$}.
    \begin{jieda}
        画图即可，$k = 1$时是函数的渐近线，可以取得. 不要忘记$k = 0$时候也是只有一个解！
    \end{jieda}
\end{enumerate}
\section*{函数性质类}
\subsection*{单调性、最值}
\begin{enumerate}
    \item 函数$y=-(x-3)\abs{x}$的严格单调增区间是\blkda{$[0,\frac{3}{2}]$}. 
    \begin{jieda}
    \ps 解析式中有部分绝对值，要考虑分类讨论，写成分段函数.
    \end{jieda}
    \item 设$f(x), g(x)$都是单调函数, 则如下四个命题中, 真命题的序号是\blkda{\ding{172},\ding{173}}: 
    \begin{dingautolist}{172}
    \item 若$f(x)$是严格增函数, $g(x)$是严格增函数, 则$f(x)+g(x)$是严格增函数; 
    \item 若$f(x)$是严格增函数, $g(x)$是严格减函数, 则$f(x)-g(x)$是严格增函数; 
    \item 若$f(x)$是严格增函数, $g(x)$是严格增函数, 则$f(x)\cdot g(x)$是严格增函数; 
    \item 若$f(x)$是严格增函数, $g(x)$是严格减函数, 则$\frac{f(x)}{g(x)}$是严格减函数. 
    \end{dingautolist}
    \item 已知函数 $f(x)=\frac{\sqrt{2021-a x}}{a-1}$ 在 $[0,1]$ 上严格单调减，那么实数 $a$ 的取值的范围是\blkda[3]{$(-\infty, 0) \cup(1,2021]$}.
    \begin{jieda}
    首先判断“定义域”，$2021-ax \geq 0$在$x \in [0, 1]$恒成立，设$g(x) = 2021-ax$，这是一条斜率不确定的线段，只需两个端点满足即可恒成立，即$g(0) \geq 0, g(1) \geq 0$，因此$a \leq 2021$，又因为$a \neq 1$，因此$a \in (-\infty, 1) \cup (1, 2021]$时$f(x)$在$[0, 1]$上有定义. \\
    当 $a<0$ 时, $y=\sqrt{2021-a x}$ 在 $[0,1]$ 上严格单调增, 且 $a-1<0$,所以 $f(x)$ 在 $[0,1]$ 上严格单调减, 符合题意, 当 $a=0$ 时, $f(x)=-\sqrt{2021}$ 并非严格增或严格减, 不符合题意, 当 $0<a<1$ 时, $y=\sqrt{2021-a x}$ 在 $[0,1]$ 上严格单调减, 且 $a-1<0$, 不符合题意, 当 $2021\geq a>1$ 时, $y=\sqrt{2021-a x}$ 在 $[0,1]$ 上严格单调减, $a-1>0$, 符合题意.\\
    综上实数 $a$ 的取值的范围是 $(-\infty, 0) \cup(1,2021]$. \\
    \ps $y = a \cdot f(x)$，即一个函数乘常系数，类型的函数的单调性在$a > 0$时与$f(x)$保持一致，在$a < 0$时与$f(x)$完全相反.
    \end{jieda}
    \item 若函数$y=\frac{a^x}{2a^2-5a+2}$在$\bR$上为严格减函数, 则$a$的取值范围是\blkda[3]{$(0,\frac{1}{2})\cup(1,2)$}. 
    \begin{jieda}
    $\begin{cases}
    a\in(0,1)\\
    2a^2-5a+2>0
    \end{cases}$
    或$\begin{cases}
    a>1\\
    2a^2-5a+2<0
    \end{cases}$, 
    解得$a\in(0,\frac{1}{2})\cup(1,2)$. 
    \end{jieda}
    \item 若函数 $f(x)=\left\{\begin{array}{l}a^x, x \geq 1 \\ \left(2-\frac{a}{2}\right) x+2, x<1\end{array}\right.$ 是 $\mathbf{R}$ 上的严格增函数, 则实数 $a$ 的取值范围是\blkda{$[\frac{8}{3}, 4)$}.
    \begin{jieda}
    $f(x)$ 是 $\mathbf{R}$ 上的严格增函数 $\Longleftrightarrow \left \{ \begin{array}{l}
        a > 1  \\
        2 - \frac{a}{2} > 0 \\
        a \geq  (2-\frac{a}{2}) + 2
    \end{array}\right.$，解得$a \in [\frac{8}{3}, 4)$
    \end{jieda}
    \item 已知函数 $f(x)=\left\{\begin{array}{l}x^{2}+2 a x, x \geq-1 \\ (2-a) x+2 a-5, x<-1\end{array}\right.$, 若 $f(x)$ 在 $\mathbf{R}$ 上是严格增函数, 则实数 $a$ 的取值范围是\blkda{$1 \leq a \leq \frac{8}{5}$} .
    \begin{jieda}
    由函数 $f(x)=\left\{\begin{array}{l}x^{2}+2 a x, x \geq-1 \\ (2-a) x+2 a-5, x<-1\end{array}\right.$ 在 $\mathbf{R}$ 上是增函数, 得 $\left\{\begin{array}{l}-a \leq-1 \\ 2-a>0 \\ 1-2 a \geq-(2-a)+2 a-5\end{array}\right.$,解得 $1 \leq a \leq \frac{8}{5}$,
    所以实数 $a$ 的取值范围是 $1 \leq a \leq \frac{8}{5}$.
    故答案为: $1 \leq a \leq \frac{8}{5}$
    \end{jieda}
    \item 设$a$为实数, 函数$f(x)=2x^2+(x-a)\abs{x-a}$. \\
    \begin{enumerate*}[(1)]
    \item 若$f(0)\geq 1$, 求$a$的取值范围; 
    \item 求$f(x)$的最小值. 
    \end{enumerate*}
    \begin{jieda}
    \begin{enumerate}[(1)]
    \item $f(0)=-a\abs{a}$，因为$g(x)=-x\abs{x}$在$\bR$上严格单调递, 且$g(-1)=1$. \\
    于是$f(0)=-a\abs{a}\geq 1=g(-1)\iff a\leq -1$.  \\
    \ps 想象一下$g(x)=-x\abs{x}$的图象.
    \item $f(x)=\begin{cases}
    x^2+2ax-a^2, & x\leq a\\
    3x^2-2ax+a^2, & x\geq a
    \end{cases}$. \\
    函数图象由两个二次函数图象的部分拼接而成，且拼接处连续，分析两个开口向上的二次函数的对称轴可知，一个为$x = -a$，一个为$\frac{a}{3}$ \\
    讨论拼接点$x = a$与两个对称轴不同的位置关系下函数的单调性，即可得到函数的最小值，在不同的拼接点条件下寻找全局的最小值即为所求. 
    \begin{enumerate}
        \item $a = 0$时，三者重合，此时最小值为$0$ 
        \item $a > 0$时，$-a < \frac{a}{3} < a$，故函数$f(x)$在第一段上先减后增，在第二段上单调增，因此最小值在第一段的顶点处取得（$x = -a$），最小值为$-2a^2$
        \item $a < 0$时，$a < \frac{a}{3} < -a$，故函数$f(x)$在第一段上单调递减，在第二段上先减后增，因此最小值在第二段顶点取得（$x = \frac{a}{3}$），最小值为$\frac{2a^2}{3}$
    \end{enumerate}
    综上，$a < 0$时为$\frac{2a^2}{3}$，$a \geq 0$时，$f(x)$的最小值为$-2a^2$. \\
    \ps 最值需要通过单调性来说明，分段函数分段求，分析各段上函数的单调性.
    另外, 本题的另一个麻烦之处在于函数中的分段点与参数$a$有关, 各段中的``对称轴''(单调区间的端点)亦与参数$a$有关. 
    \end{enumerate}
    \end{jieda}
    \item 设函数$f(x)=\frac{a(a^{2x}-1)}{a^x(a^2-1)}$($a>0$且$a\neq 1$), 证明: $y=f(x)$是$\bR$上为严格增函数. 
    \begin{jieda}
    $f(x)=\frac{a}{a^2-1}\cdot(a^x-a^{-x})$. \\
    任取实数集内的$x_1, x_2, x_1 < x_2$ \\
    $f(x_2)-f(x_1)=\frac{a}{a^2-1}\left(a^{x_2}-a^{x_1} - (a^{-x_2} - a^{-x_1})\right)$. 
    \begin{enumerate}
        \item 当$0 < a < 1$时,$a^{x_2}-a^{x_1} < 0$，$a^{-x_2} - a^{-x_1} > 0$，$\frac{a}{a^2-1} < 0$，故$f(x_2)-f(x_1) > 0$
        \item 当$a > 1$时,$a^{x_2}-a^{x_1} > 0$，$a^{-x_2} - a^{-x_1} < 0$，$\frac{a}{a^2-1} > 0$，故$f(x_2)-f(x_1) > 0$
    \end{enumerate}
    因此$f(x)$是$\bR$上为严格增函数. \\
    \ps $f(x)$并非恒正或者恒负，因此不应使用作商法.
    \end{jieda}
    \item 已知$(a+1)^{-\frac{3}{5}}<(3-2a)^{-\frac{3}{5}}$, 求实数$a$的取值范围. 
    \begin{jieda}
    \begin{enumerate}
        \item 方法1 \\
        构建函数$f(x) = x^{-\frac{3}{5}}$，原不等式转化为$f(a+1) < f(3-2a)$ \\
        易知函数$f(x) = x^{-\frac{3}{5}}$有两个单调减区间，奇函数，在$x > 0$时$f(x) > 0$，在$x < 0$时$f(x) < 0$ 
        \begin{enumerate}
            \item 若$(a+1)(3-2a) > 0$，即$a \in (-1, \frac{3}{2})$，则$a+1$和$3-2a$位于同一个单调区间，此时根据单调性，$f(a+1) < f(3-2a)$被转化为$a+1 > 3-2a$。故$a \in (\frac{2}{3}, \frac{3}{2})$
            \item 若$a+1< 0 < 3-2a$，即$a \in (-\infty, -1)$，则$a+1$和$3-2a$位于不同单调区间，并且函数的性质可知满足$f(a+1) < f(3-2a)$
        \end{enumerate}
        综上，$a \in (-\infty, -1) \cup (\frac{2}{3}, \frac{3}{2})$
        \item 方法2 \\
        构建函数$f(x) = x^{\frac{3}{5}}$，原不等式转化为$f(\frac{1}{a+1}) < f(\frac{1}{3-2a})$ \\
        易知函数$f(x) = x^{\frac{3}{5}}$为$R$上单调递增的奇函数，因此$f(\frac{1}{a+1}) < f(\frac{1}{3-2a})$即$\frac{1}{a+1} < \frac{1}{3-2a}$ \\
        移项通分有$\frac{2-3a}{(a+1)(3-2a)} < 0$，解集为$a \in (-\infty, -1) \cup (\frac{2}{3}, \frac{3}{2})$
    \end{enumerate}
    \end{jieda}
    \item 函数$y=3^{\sqrt{x^2-5x+4}}$的递减区间是\blkda{$(-\infty,1]$}, 值域是\blkda{$[1,+\infty)$}. 
    \begin{jieda}
    定义域$x\in(-\infty,1]\cup[4,+\infty)$. \\
    根据复合函数的单调区间的原则``同增异减''可得递减区间是$(-\infty,1]$. \\
    $\sqrt{x^2-5x+4}\in[0,+\infty)$, 故$y\in[1,+\infty)$. 
    \end{jieda}
    \item 函数$y = \frac{1}{x^2+2x-4}$的严格增区间是\blkda[5]{$(-\infty,-1-\sqrt{5})$和$(-1-\sqrt{5}, -1]$}. 
    \begin{jieda}
        定义域先行：$x \in \{x | x \neq -1\pm \sqrt{5}\}$ \\
        外层函数为反比例函数，有两个严格减区间，故要寻找整体的严格增区间，应寻找内层函数的严格减区间即$(-\infty,-1]$ \\
        综合两者，因此是两个单调区间：$(-\infty,-1-\sqrt{5})$和$(-1-\sqrt{5}, -1]$ \\
        \ps 单调区间之间不能用并集符号连接！
    \end{jieda}
    \item $f(x)   = x^2$, 设函数$g(x)=-q\cdot f(f(x))+(2q-1)\cdot f(x)+1$, 问是否存在实数$q(q<0)$, 
    使得$g(x)$在区间$(-\infty,-4]$上是严格减函数, 且在$(-4,0)$上是严格增函数? 若存在, 请求出$q$; 若不存在, 请说明理由.
    \begin{jieda}
    $g(x)=-qx^4+(2q-1)x^2+1$, 设$g(x) = h(f(x))$，$h(u) = -qu^2 + (2q-1)u + 1, u = f(x) = x^2$ \\
    内层函数$f(x)$在$(-\infty, 0)$严格减，在$[0, +\infty)$严格增，值域为$[0， +\infty)$ \\
    外层函数$h(u)$在$(-\infty, \frac{2q-1}{2q})$严格减，在$[\frac{2q-1}{2q}, +\infty)$严格增 \\
    由$q < 0$知$\frac{2q-1}{2q} > 0$，故当$x \in (-\infty, -\sqrt{\frac{2q-1}{2q}}), f(x)$严格减，$u \geq \frac{2q-1}{2q}, h(u)$严格增，故$g(x)$严格减; 当$x \in (-\sqrt{\frac{2q-1}{2q}}, 0), f(x)$严格减，$0 < u < \frac{2q-1}{2q}, h(u)$严格减，故$g(x)$严格增. 因此，当$-\sqrt{\frac{2q-1}{2q}} = -4$，即$q=-\frac{1}{30}$时符合题目要求. 
    \end{jieda}
    \item 定义在 $(0,+\infty)$ 上的函数 $f(x)$ 满足 $\frac{x_{2} f\left(x_{1}\right)-x_{1} f\left(x_{2}\right)}{x_{1}-x_{2}}<0$, 且 $f(3)=9$, 则使 $f(x)>3 x$ 成立的 $x$ 的取值范围是\blkda{$(0,3)$} .
    \begin{jieda}
    由 $\frac{x_{2} f\left(x_{1}\right)-x_{1} f\left(x_{2}\right)}{x_{1}-x_{2}}<0$, 且 $x_{1}, x_{2} \in(0,+\infty)$, 则两边同时除以 $x_{1} x_{2}$ 可得 $\frac{\frac{f\left(x_{1}\right)}{x_{1}}-\frac{f\left(x_{2}\right)}{x_{2}}}{x_{1}-x_{2}}<0$,
    令 $g(x)=\frac{f(x)}{x}$, 则原不等式为 $\frac{g\left(x_{1}\right)-g\left(x_{2}\right)}{x_{1}-x_{2}}<0$, 因此函数 $g(x)$ 在 $(0,+\infty)$ 上严格减,由 $f(x)>3 x$, 得 $\frac{f(x)}{x}>3$, 又 $g(3)=\frac{f(3)}{3}=3$, 于是 $g(x)>g(3)$, 解得 $0<x<3$,所以使 $f(x)>3 x$ 成立的 $x$ 的取值范围是 $(0,3)$. \\
    \ps $\frac{f(x_1) - f(x_2)}{x_1 - x_2}$的几何意义是割线的斜率.
    \end{jieda}
\end{enumerate}
\subsection*{奇偶性、对称性}
\begin{enumerate}
    \item 函数$f(x)=x\abs{x+a}+b$是奇函数的充要条件是\dotfill(\kh{B})
    \xx{$ab=0$}{$a^2+b^2=0$}{$a=b$}{$a+b=0$}
    \begin{jieda}
    奇函数若在$x=0$处有定义, 则此处的函数值为零, 即$f(0)=0$, 于是$b=0$. \\
    故$f(x)=x\abs{x+a}$. 于是由$f(x)$为奇函数，不难得到$a=0$, 比如$f(-a)+f(a)=0$. 
    \end{jieda}
    \item 若函数$f(x)=\frac{a-\me^x}{1+a\cdot\me^x}$($a$为常数)在定义域上为奇函数, 则$a$的值为\blkda{$\pm 1$}. 
    \begin{jieda}
    首选解法：$f(0)=0$或无意义，即$\frac{a-1}{1+a} = 0$或$\frac{a-1}{1+a}$无意义，故$a = \pm 1$\\
    也可以采用就极限思想：让函数在$\pm \infty$时的极限互为相反数求解：$x \rightarrow +\infty, f(x) \rightarrow -\frac{1}{a}, x \rightarrow -\infty, f(x) \rightarrow a$，由$a = \frac{1}{a}$得$a = \pm 1$. \\
    把握定义求解有：$f(x) = \frac{a-\me^x}{1+a\cdot\me^x}, f(-x) = \frac{a-\me^{-x}}{1+a\cdot\me^{-x}} = \frac{a\me^{x}-1}{a+\me^{x}}$ \\
    由$f(x) = -f(-x)$知$\frac{a-\me^x}{1+a\cdot\me^x} = \frac{1-a\me^{x}}{a+\me^{x}}$，交叉相乘得到$a^2 - e^{2x} = 1-a^2e^{2x}$，移项整理得$(a^2-1) = (1-a^2)\cdot(e^{2x})$，因为此式恒成立，所以必有$a = \pm 1$
    \end{jieda}
    \item 已知函数$f(x)=x^5+ax^3+bx-8$, 且$f(-2)=10$, 则$f(2)$=\blkda{$-26$}
    \begin{jieda}
    $f(x)+f(-x)=-16$，函数关于$(0, -8)$中心对称。\\
    \ps 设$g(x) = f(x)+8$，则$g(x)$为奇函数. (多项式函数是奇函数等价于偶次方的系数为零, 是偶函数等价于奇次方的系数为零.)
    \end{jieda}
    \item 设$f(x)$和$g(x)$分别是$\mathbb{R}$上的偶函数和奇函数, 则下列结论恒成立的是\dotfill(\kh{A})
    \xx{$f(x)+\abs{g(x)}$是偶函数}{$f(x)-\abs{g(x)}$是奇函数}{$\abs{f(x)}+g(x)$是偶函数}{$\abs{f(x)}-g(x)$是奇函数}
    \begin{jieda}
    $g(x)$为奇函数，则$|g(x)|$为偶函数，偶函数$+$偶函数还是偶函数，故$A, B$都是偶函数，$A$正确，$B$错误。$f(x)$是偶函数，$|f(x)|$还是偶函数。奇函数$+$偶函数的结果是不确定的（如果两个函数都不是“既是奇函数又是偶函数”的话，结果必然非奇非偶函数）。故$C,D$都错。\\
    \ps 任意一个关于定义域关于原点对称的函数都可以写成一个奇函数加上一个偶函数（包含其中之一为既奇又偶的情况）：$h(x) = \frac{h(x)+h(-x)}{2} + \frac{h(x) - h(-x)}{2}$
    \end{jieda}
    \item 已知$m$是正整数，幂函数$f(x) = x^{1+m-m^2}$在$(0, +\infty)$上时严格增函数.
\begin{enumerate}
    \item 求幂函数$f(x)$的解析式
    \item 设$\varphi(x) = \frac{a[f(x)]^2+1}{bf(x)+c}$（$a, b, c$为整数），若$\varphi(x)$是奇函数，又$\varphi(1) = 2, \varphi(2) < 3$，求$a, b, c$的值.
\end{enumerate}
\begin{jieda}
    \begin{enumerate*}
        \item 根据$f(x)$在$(0, +\infty)$上时严格增函数知$1+m-m^2>0$，又因为$m$是正整数，所以$m = 1$,故$f(x) = x$ \\
        \item $\varphi(x) = \frac{ax^2+1}{bx+c}$. 因为$\varphi(x)$是奇函数，故$0$不在定义域内，否则$\frac{1}{c}$必不为0，因此$c = 0$，经检验，此时$\varphi(x)$为奇函数. $\varphi(1) = \frac{a+1}{b} = 2$，即$b = \frac{a+1}{2}$；$\varphi(2) = \frac{4a+1}{2b} = \frac{4a+1}{a+1} = 4 - \frac{3}{a+1} < 3$，故$a \in \{0, 1\}$，又因为$b$是整数，故$a = 1$\\
        综上，$a = 1, b = 1, c = 0$
    \end{enumerate*}
    \end{jieda}    
    \item 已知函数$y=\frac{1}{4-2^x}$的图像关于点$P$对称, 则点$P$的坐标为\dotfill(\kh{C})
    \xx{$(2,\frac{1}{2})$}
    {$(2,\frac{1}{4})$}
    {$(2,\frac{1}{8})$}
    {$(0,0)$}
    \begin{jieda}
    设$P(a,b)$, 因为$x\neq 2$, 故$a=2$. \\
    $2b=f(1)+f(3)$, 于是$b=\frac{1}{8}$. \\
    \ps $f(x)$的对称中心是$(a, b)$或对称轴是$x = a$的必要条件均为$f(x)$的定义域关于$a$对称.
    \end{jieda}
    \item 已知函数$f(x) = -x^3-x$，若$f(t^2+2t+k)+f(-2t^2+2t-5)>0$对任意的$t \in [0, 5]$都成立，则$k$的取值范围是\blkda{$(-\infty, 1)$}.
    \begin{jieda}
        $f(x)$为奇函数，在$\mathbf{R}$上单调递减，根据对称性可以将函数值上的不等关系转化为变量的不等关系，即$(t^2+2t+k)+(-2t^2+2t-5)<0$在$t \in [0, 5]$恒成立.\\
        参变分离有$-t^2+4t-5 < -k$在$t \in [0, 5]$恒成立，因此有$k < 1$ \\
        \ps $f(x)$关于$(a, b)$对称，则$f(x_1) + f(x_2) = 2b \Leftrightarrow x_1 + x_2 = 2a$ \\
        $f(x)$关于$(a, b)$对称，且在定义域内单调增，则$f(x_1) + f(x_2) < 2b \Leftrightarrow x_1 + x_2 < 2a$ \\
        $f(x)$关于$(a, b)$对称，且在定义域内单调增，则$f(x_1) + f(x_2) > 2b \Leftrightarrow x_1 + x_2 > 2a$ \\
        $f(x)$关于$(a, b)$对称，且在定义域内单调减，则$f(x_1) + f(x_2) > 2b \Leftrightarrow x_1 + x_2 < 2a$ \\
        $f(x)$关于$(a, b)$对称，且在定义域内单调减，则$f(x_1) + f(x_2) < 2b \Leftrightarrow x_1 + x_2 > 2a$ \\
        以上结论可以通过代数推导，更推荐结合图像理解.
    \end{jieda}
    \item 已知函数$f(x) = 2^x - (\frac{1}{2})^x+3$，且$f(m^2)+f(-5m+4)<6$，则实数$m$的取值范围是\blkda{$(1, 4)$}.
    \begin{jieda}
        直接应用上面结论即可，也可以设$g(x) = f(x) - 3$，将$f(m^2)+f(-5m+4)<6$转化为$g(m^2)+g(-5m+4)<0$后按照更熟悉的奇函数形式处理.
    \end{jieda}
    \item 已知函数$y = f(3x+1)$为偶函数，且在$[0, +\infty)$上为严格增函数，若$f(x) < f(2x+1)$，则$x$的取值范围是\blkda[5]{$(-\infty, -1) \cup (\frac{1}{3}, +\infty)$}.
    \begin{jieda}
        函数$y = f(3x+1)$为偶函数，且在$[0, +\infty)$上为严格增函数即$f(x)$关于$x = 1$对称，且“开口向上”，因此离对称轴$x = 1$近的变量的函数值小，由此$f(x) < f(2x+1)$可以转化为$|x-1| < |2x|$，两侧平方有$(x-1)^2 < 4x^2$，解得$x \in (-\infty, -1) \cup (\frac{1}{3}, +\infty)$
    \end{jieda}
    \item 已知函数$f(x)$在定义域$[2-a, 3]$上是偶函数，在$[0, 3]$上严格单调递减，并且$f(-m^2-\frac{a}{5}) > f(-m^2+2m-2)$，则$m$的取值范围是\blkda[3]{$[1-\sqrt{2}, \frac{1}{2}]$}.
    \begin{jieda}
        首先根据定义域的对称性可以求得$a = 5$，进一步根据对称性可以判断出函数“开口向下”，离$y$轴越近函数值越大，再结合定义域的要求，则$f(-m^2-\frac{a}{5}) > f(-m^2+2m-2) \Leftrightarrow \left \{ \begin{array}{l}
             |-m^2-1| < |-m^2+2m-2|  \\
             -3 \leq -m^2-1 \leq 3  \\
             -3 \leq -m^2+2m-2 \leq 3 
        \end{array}\right.$第一个配方后可以直接去绝对值，最终得到$m \in [1-\sqrt{2}, \frac{1}{2}]$
    \end{jieda}
    \item 已知$f(x) = \frac{(2^x+1)^2}{2^x \cdot x}+1$在$[-2018, -1]\cup[1, 2018]$上最大值为$M$,最小值为$m$,则$M+m = $\blkda{$2$}.
    \begin{jieda}
    设$g(x) = \frac{(2^x+1)^2}{2^x \cdot x}$，则$g(x) = \frac{2^x+\frac{1}{2^x}+2}{x} = \frac{2^x+2^{-x}+2}{x}$，故$g(-x) = -g(x)$，即$g(x)$为奇函数，故$M+m = 2$ \\
    \ps 涉及最大值与最小值和的题目一般可以优先考虑对称性.
    \end{jieda}
    \item 已知函数$f(x) = \frac{3^{x+1}+1}{3^x+1}$在$[-2021, 2021]$上的最大值和最小值分别为$M, m$，则$M+m = $\blkda{$4$}.
    \begin{jieda}
        $g(x) = f(x) - 2 = \frac{3^x - 1}{3^x+1}$为奇函数，故$f(x)$关于$(0, 2)$中心对称，$M+m = 4$ \\
        \ps 本题如果没有发现$f(x) - 2$是奇函数，利用单调性也是可以解决的，计算量相对大. \\
        \ps 其实$y=\frac{a^x-1}{a^x+1}$, $y=\frac{a^x+1}{a^x-1}$都是奇函数.
    \end{jieda}
    \item 已知函数$f(x) = x^3 + x (x \in \mathbf{R})$，函数$g(x)$满足$g(x)+g(2-x) = 0(x \in \mathbf{R})$，若函数$h(x) = f(x-1)-g(x)$恰有2019个零点，则所有这些零点之和为\blkda{$2019$}.
    \begin{jieda}
    易知$f(x)$为奇函数，故$f(x-1)$关于$(1, 0)$中心对称。又由$g(x)+g(2-x) = 0(x \in \mathbf{R})$可知$g(x)$关于$(1, 0)$中心对称。因此$h(x)$关于$(1, 0)$中心对称，因此其零点必然关于$1$对称（奇数个零点代表有一个零点为1），故所有零点之和为零点个数，2019.
    \end{jieda}
    \item 函数$f(x)$为定义在$(-\infty, 0) \cup (0, +\infty)$上的奇函数，且$f(2) = 1$，对于任意$x_1, x_2 \in (0, +\infty), x_1 \neq x_2$，都有$\frac{x_1f(x_1)-x_2f(x_2)}{x_1 - x_2} > 0$成立，则$f(x) \leq \frac{2}{x}$的解集为\blkda[3]{$(-\infty, -2] \cup (0, 2]$}.
    \begin{jieda}
        设$g(x) = xf(x)$，则对于任意$x_1, x_2 \in (0, +\infty), x_1 \neq x_2$，都有$\frac{g(x_1)-g(x_2)}{x_1 - x_2} > 0$成立，即$g(x)$在$(0, +\infty)$严格增；又因为$f(x)$为奇函数，所以$g(x)$为偶函数. \\
        不等式$f(x) \leq \frac{2}{x}$转化为$x > 0, g(x) \leq 2$或$x < 0, g(x) \geq 2$，解得$x \in (-\infty, -2] \cup (0, 2]$ \\
        \ps 本题转化不等式的时候一定要注意分类讨论
    \end{jieda}
    \item 已知二次函数 $f(x)=a x^2+b x(a \neq 0)$ 满足: $f(5-x)=f(x-3)$ 且方程 $f(x)=x$ 有等根. 
    \begin{enumerate}
    \item 求 $f(x)$ 的解析式;
    \item 是否存在实数 $m, n(m<n)$, 使得 $f(x)$ 的定义域为 $[m, n]$, 值域为 $[k m, k n](k>1)$.
    \end{enumerate}
    \begin{jieda}
    \begin{enumerate*}
    \item $\because f(5-x)=f(x-3)$
    $\therefore f(x)$ 的对称轴为 $x=\frac{5-x+x-3}{2}=1$
    即 $-\frac{b}{2 a}=1$, 即 $b=-2 a$
    $\because f(x)=x$ 有两个相等的实根
    $\therefore a x^2+b x=x$
    即 $a x^2+(b-1) x=0$, 有两个相等的实根 0
    $\therefore x_1+x_2=-\frac{b-1}{a}=0$
    $\therefore b=1, a=-\frac{1}{2}$
    $\therefore f(x)=-\frac{1}{2} x^2+x$
    综上所述, 结论是: $f(x)$ 的解析式为
    $f(x)=-\frac{1}{2} x^2+x$ \\
    \item $f(x)=-\frac{1}{2} x^2+x=-\frac{1}{2} \times(x-1)^2+\frac{1}{2} \leq \frac{1}{2}$
    $\therefore k n \leq \frac{1}{2}, \therefore m<n<\frac{1}{2k}$
    又 $\because$ 函数的对称轴为 $x=1$
    $\therefore f(x)$ 在 $[m, n]$ 单调递增则有 $f(m)=k m, f(n)=k n$
    故$m, n$是$-\frac{1}{2}x^2 + x = kx$的两根，故$m = 2(1-k), n = 0$ \\
    \ps 根据值域判断函数在定义域内的单调性是简化，解决问题的突破口。\\
    常规解决方法是分析单调性，根据最值列方程组.
    \end{enumerate*}
    \end{jieda}
    \item 函数$f(x) = \frac{x}{ax+b}$（$a, b$是非零常数），满足$f(2) = 1$，且方程$f(x) = x$仅有一个解.
    \begin{enumerate}
        \item 求$a, b$的值
        \item 是否存在常数$m$，使得对于定义域内任意的$x$，$f(x) + f(m-x) = 4$恒成立
        \item 在平面直角坐标系中，求定点$A(-3, 1)$到此函数上任意一点$P$的距离$|AP|$的最小值.
    \end{enumerate}
    \begin{jieda}
        \begin{enumerate*}
            \item 由$f(2) = 1$，$\frac{2}{2a+b} = 1$，故$2a+b = 2$；$f(x) = x$即$\frac{x}{ax+b} = x$，即$(\frac{1}{ax+b}-1)x = 0$，容易判断其有两个解$x_1  = 0, x_2 = \frac{1-b}{a}$.要符合题意只可能是$b = 1$，两根重合.进一步推得$a = \frac{1}{2}$ \\
            \item $f(x) = \frac{2x}{x+2}$，为反比例型函数，其对称中心为$(-2, 2)$，因此$f(x)$满足$f(x) + f(-4-x) = 4$. 由此可知$m = -4$ \\
            \item 设$P(x_0, \frac{2x_0}{x_0+2})$，故$|AP| = \sqrt{(x_0+3)^2+(\frac{2x_0}{x_0+2}-1)^2}   = \sqrt{(x_0+3)^2+(1-\frac{4}{x_0+2})^2}$. 设$t = x_0+2 (t \neq 0)$，则$|AP| = \sqrt{(t+1)^2+(1-\frac{4}{t})^2} = \sqrt{(t-\frac{4}{t})^2+2(t-\frac{4}{t})+10}$，设$k = t - \frac{4}{t}, k \in \mathbf{R}$，$|AP| = \sqrt{k^2+2k+10} = \sqrt{(k+1)^2+9} \geq 3$（当$x = \frac{-5\pm\sqrt{17}}{2}$时取得最小值）
        \end{enumerate*}
    \end{jieda}
    \item 函数$y=f(x-a)$与函数$y=f(-x+a)$的图像关于\blkda{$x=a$} 对称. 
    \begin{jieda}
        $y=f(x-a)$与$y=f(-x+a)$是两个不同的函数！这里是互对称场景.
    \end{jieda}
\end{enumerate}
\subsection*{周期性}
\begin{enumerate}
    \item 定义在$\bR$上的函数$f(x)$的图像关于点$(a,b),(c,b)$都对称($a\neq c$), 则\dotfill(\kh{B})
    \xx{$f(x)$是以$\abs{a-c}$为周期的周期函数}
    {$f(x)$是以$2\abs{a-c}$为周期的周期函数}
    {$f(x)$是以$\frac{1}{2}\abs{a-c}$为周期的周期函数}
    {$f(x)$不是周期函数}
    \begin{jieda}
    两个横坐标不同位置的对称性会产生周期性
    \begin{enumerate}
        \item 两个对称轴的横坐标之差的二倍是函数的周期：$f(x) = f(2a-x), f(x) = f(2b - x) \Leftrightarrow 2|a-b|$是$f(x)$的一个周期。
        \item 两个纵坐标相同的对称中心的横坐标之差的二倍是函数的周期：$f(x) + f(2a-x) = 2m, f(x) + f(2b - x) = 2m \Leftrightarrow 2|a-b|$是$f(x)$的一个周期。
        \item 一个对称轴和一个对称中心的横坐标之差的四倍是函数的周期：$f(x) = f(2a-x), f(x) + f(2b - x) = 2m \Leftrightarrow 4|a-b|$是$f(x)$的一个周期。
    \end{enumerate}
    \end{jieda}
    \item 设$f(x)$是定义在$\bR$上的奇函数, 且对于任意的$x\zr$, $f(1+x)-f(1-x)=0$恒成立; 当$x\in[0,1]$时, $f(x)=2x$. 
    若方程$f(x)=ax$恰好有5个不同的解, 则实数$a$的取值范围为\blkda[2.5]{$(-\frac{2}{3},-\frac{2}{7})\cup\{\frac{2}{5}\}$}. 
    \begin{jieda}
    $f(x)$图像关于原点对称, 关于直线$x=1$对称, 可知周期$T=4$. \\
    $y = ax$的图像是过原点的直线，随着$a$的变化，其斜率变化. 绘制草图后可以判断要分为$a>0$与$a<0$两种情形, 
    不难求得$a=\frac{2}{5}$或$a\in(-\frac{2}{3},-\frac{2}{7})$. 
    \end{jieda}
    \item $f(x)$是定义在$\mathbf{R}$上的以3为周期的奇函数，且$f(2) = 0$，则方程$f(x) = 0$在区间$(0, 6)$内的解的个数的最小值是\blkda{$7$}. 
    \begin{jieda}
        $f(x)$为$\mathbf{R}$上奇函数，故$f(0) = 0$，又因为$f(2) = 0$，由周期性$f(-1) = 0$，由奇函数性质进一步可以推得$f(-2) = 0$，故由周期性可知$f(n) = 0 (n \in \mathbf{Z})$ \\
        由奇函数和周期性可知$f(1.5) = f(-1,5) = 0$，进一步可知$f(4.5) = 0$ \\
        \ps 对于周期为$T$的奇函数，必有$f(\frac{T}{2}) = f(-\frac{T}{2}) = 0$
    \end{jieda}
    \item 已知函数$f(x)\neq -1$, 且对定义域内任意的$x$总有关系$(f(x+\pi)+1)(f(x)+1)=2$, 那么下列结论中正确的是\dotfill(\kh{B})
    \xx{$f(x)$是周期为$\pi$的周期函数}
    {$f(x)$是周期为$2\pi$的周期函数}
    {$f(x)$是周期为$\frac{\pi}{2}$的周期函数}
    {$f(x)$不是周期函数}
    \begin{jieda}
    由条件可得$(f(x+2\pi)+1)(f(x+\pi)+1)=2$, 于是$f(x+2\pi)=f(x)$. \\
    \ps $f(x) + f(x+t) = c\Leftrightarrow 2t$是$f(x)$的一个周期。\\
    进阶版：$f(x+a) + f(x+b) = c\Leftrightarrow 2|a-b|$是$f(x)$的一个周期。 \\
    $f(x) \cdot f(x+t) = c (c \neq 0)\Leftrightarrow 2t$是$f(x)$的一个周期。\\
    进阶版：$f(x+a) \cdot f(x+b) = c (c \neq 0)\Leftrightarrow 2|a-b|$是$f(x)$的一个周期。
    \end{jieda}
    \item 已知$f(x)$是$\bR$上的函数, 且$f(x+2)(1-f(x))=1+f(x)$. 
    \begin{enumerate*}[(1)]
    \item 求证: $f(x)$是周期函数;
    \item 若$f(4)=-\sqrt{3}$, 求$f(2020)$. 
    \end{enumerate*} 
    \begin{jieda}
    \begin{enumerate*}[(1)]
    \item 若$f(x)$为常数函数, 则设为$f(x)=t$, 于是$t(1-t)=1+t$, 无解, 故$f(x)$不是常函数. \\
    (i) 若$f(x)=1$有解, 设为$x_0$, 则$f(x_0+2)(1-f(x_0))=1+f(x_0)$, 故$f(x_0)=-1$. 矛盾. \\
    若$f(x)=0$有解, 设为$x_0$, 则$f(x_0+2)(1-f(x_0))=1+f(x_0)$, 故$f(x_0+2)=1$. 显然不可能(见上一行). \\
    (ii) 由(i)知$f(x)=1$无解, 且$f(x)=0$无解, \\
    则$f(x+4)=\frac{1+f(x+2)}{1-f(x+2)}=\frac{1+\frac{1+f(x)}{1-f(x)}}{1-\frac{1+f(x)}{1-f(x)}}=-\frac{1}{f(x)}$, 故$f(x+8) = -\frac{1}{f(x+4)} = f(x)$. \\
    综上, $f(x)$是周期为8的周期函数. \\
    \item $f(2020)=f(4)=-\sqrt{3}$. 
    \end{enumerate*}
    \end{jieda}
    \item 已知$f(x)$是定义域为$(-\infty, +\infty)$的奇函数，满足$f(1-x) = f(1+x)$。若$f(1) = 2$，则$f(1)+f(2)+f(3)+...+f(50) = $\dotfill(\kh{C})
    \xx{-50}
    {0}
    {2}
    {50}
    \begin{jieda}
        因为$f(x)$为奇函数，且$f(1-x) = f(1+x)$，故$f(x)$具有周期性，最小正周期为4，即$f(x) = f(x+4)$. 由题意可知$f(0) = 0, f(1) = 2$, 可以推出$f(2) = 0, f(3) = f(-1) = -2$，由此可知$f(4k)+f(4k+1)+f(4k+2)+f(4k+3)=0, k \in \mathbf{Z}$。故原式$= f(1) + f(2) = 2$ \\
        \ps 具有周期性的求和问题，先求一个周期内的和.
    \end{jieda}
\end{enumerate}
\section*{基本初等函数}
\subsection*{幂函数}
\begin{enumerate}
    \item 已知幂函数 $y=\left(m^2-m-1\right) x^{\frac{m}{3}}$ 的图像关于 $y$ 轴对称, 则 $m=$ \blkda{$2$}.
    \item 如图曲线是幂函数$y=x^n$在第一象限内的图像, 已知$n$取$\pm2,\pm\frac{1}{2}$四个值, 则曲线$c_1,c_2,c_3,c_4$对应的值依次为\dotfill(\kh{B})
    \xx{$-2,-\frac{1}{2},\frac{1}{2},2$}
    {$2,\frac{1}{2},-\frac{1}{2},-2$}
    {$-\frac{1}{2},-2,2,\frac{1}{2}$}
    {$2,\frac{1}{2},-2,-\frac{1}{2}$}
    \begin{minipage}{\linewidth}
    \centering
    
\includegraphics[totalheight=3cm]{./bodyxb/src/main42chap21mihanshu01.png}
    \end{minipage}
    
    
    \item 如图是函数$y=x^{\frac{m}{n}}$($m,n\zn$且互质)的图像, 则\dotfill(\kh{C})
    \xx{$m,n$是奇数, 且$\frac{m}{n}<1$}
    {$m$是偶数, $n$是奇数, 且$\frac{m}{n}>1$}
    {$m$是偶数, $n$是奇数, 且$\frac{m}{n}<1$}
    {$m$是奇数, $n$是偶数, 且$\frac{m}{n}>1$}
    \begin{minipage}{\linewidth}
    \centering
    \includegraphics[totalheight=3cm]{./bodyxb/src/main42chap21mihanshu02.png}
    \end{minipage}
    \item 下列命题中, 真命题的是\dotfill(\kh{C})
    \xx{幂函数中不存在既不是奇函数, 又不是偶函数的函数;}
    {如果一个幂函数不是偶函数, 那么它一定是奇函数;}
    {图像不经过$(-1,1)$的幂函数, 一定不是偶函数;}
    {如果两个幂函数有三个公共点, 那么这两个函数一定相同.}
    \item 记$\prod\limits_{k=1}^na_k=a_1\cdot a_2\cdot a_3\cdot\cdots\cdot a_n$, 
    设$f_n(x)=x^{\frac{n}{3}}, g_n(x)=x^{\frac{n}{5}}, h_n(x)=x^{\frac{n}{2}}$, 
    下列论断中, 正确论断的序号为\blkda{\ding{172}}.
    \begin{dingautolist}{172}
    \item 使函数$F(x)=\prod\limits_{k=1}^nf_k(x)$是偶函数的最小正整数$n$的值等于3;
    \item 使函数$G(x)=\prod\limits_{k=1}^ng_k(x)$是偶函数的最小正整数$n$的值等于4;
    \item 使函数$H(x)=\prod\limits_{k=1}^nh_k(x)$是奇函数的最小正整数$n$的值等于3.
    \end{dingautolist}
    \begin{jieda}
    积函数的奇偶性. 
    \begin{dingautolist}{172}
    \item $f_n(x)$为奇函数$\iff n$为奇数, $f_n(x)$为偶函数$\iff n$为偶数, 故真. 
    \item $g_n(x)$为奇函数$\iff n$为奇数, $g_n(x)$为偶函数$\iff n$为偶数, 故假. $n_{\min}=3$. 
    \item $h_1(x)$的定义域为$[0,+\infty)$, 故非奇非偶. 故假. 
    \end{dingautolist}
    \end{jieda}
    \item 已知不等式$x^2+(a-2)x+3a>0(*)$
    \begin{enumerate}
        \item 若不等式$(*)$对任意的$x \in [-1, 1]$恒成立，求实数$a$的取值范围
        \item 若不等式$(*)$对任意的$a \in [-1, 1]$恒成立，求实数$x$的取值范围
    \end{enumerate}
    \begin{jieda}
        \begin{enumerate}
            \item 
            \begin{enumerate}
                \item 通解法 \\
                $f(x) = x^2+(a-2)x+3a$的对称轴为$x = 1-\frac{a}{2}$
                \begin{dingautolist}{172}
                    \item 若$1-\frac{a}{2} < -1$，即$a > 4$，此时$f(x)$在$[-1, 1]$上单调递增，因此$f(-1) > 0$，即$1+2-a+3a > 0$，解得$a > -\frac{3}{2}$，结合对称轴有$a > 4$
                    \item 若$1-\frac{a}{2} > 1$，即$a < 0$，此时$f(x)$在$[-1, 1]$上单调递减，因此$f(1) > 0$，即$1+a-2+3a > 0$，解得$a > \frac{1}{4}$，此种情况无解
                    \item 若$1-\frac{a}{2} \in [-1, 1]$，即$a \in [0, 4]$，此时$f(x)$在$[-1, 1]$内最小值在顶点处取，即$f(1-\frac{a}{2}) > 0$，整理有$a^2-16a+4<0$，解得$a \in (8-2\sqrt{15}, 8+2\sqrt{15})$，结合对称轴有$a \in (8-2\sqrt{15}, 4]$
                \end{dingautolist}
                综上，$a \in (8-2\sqrt{15}, +\infty)$
                \item 参变分离 \\
                $(x+3)a-2x+x^2 > 0$在$x \in [-1, 1]$恒成立，即$a > \frac{2x-x^2}{x+3}$在$x \in [-1, 1]$恒成立.\\
                换元$t = x+3$， $t \in [2, 4]$，即$a > \frac{-t^2+8t-15}{t}$在$[2, 4]$恒成立.进一步变形有$a > -(t+\frac{15}{t})+8$在$[2, 4]$恒成立.根据对勾函数性质可知$-(t+\frac{15}{t})+8 \leq 8 - 2\sqrt{15}$（$t = \sqrt{15}$，即$x = \sqrt{15}-3$时取等号）,故$a \in (8-2\sqrt{15}, +\infty)$
            \end{enumerate}
            \item 将$a$做为主元，即$(x+3)a-2x+x^2 > 0$在$a \in [-1, 1]$恒成立，设$f(a) = (x+3)a-2x+x^2$，故只需$f(1) > 0, f(-1) > 0$即可.解得$x \in (-\infty, \frac{3-\sqrt{21}}{2})\cup(\frac{3+\sqrt{21}}{2}, +\infty)$
        \end{enumerate}
        \ps 恒成立问题，首先考虑主元应该是谁，然后考虑是否要参变分离.
    \end{jieda}
\end{enumerate}
\subsection*{指数函数、指数方程}
\begin{enumerate}
    \item 如图, 指数函数$y=a^x$, $y=b^x$, $y=c^x$, $y=d^x$, 则$a,b,c,d$与1的大小关系为\dotfill(\kh{B})
    \xx{$a<b<1<c<d$}
    {$b<a<1<d<c$}
    {$1<a<b<c<d$}
    {$a<b<1<d<c$}
    \begin{minipage}{\linewidth}
    \centering
    \includegraphics[totalheight=3cm]{./bodyxb/src/main42chap21zhishuhanshu01.png}
    \end{minipage}
    \ps 取横坐标为1，判断纵坐标即可.
    \item 函数$y = x^2$与$y = 2^x$的交点个数为\blkda{$3$}. 
    \begin{jieda}
        指数函数增长得很快（$a > 1$），当其从上方穿越了一个多项式函数后一定会从下方再次穿越这个多项式函数，$0 < a < 1$则是对称的情况.
    \end{jieda}
    \item 已知函数$f(x)=4+a^{x-1}$的图像恒过定点$P$, 则$P$点坐标是\blkda{$(1,5)$}. 
    \begin{jieda}
    $a^0=1$!!!
    \end{jieda}
    \item 已知函数$f(x)=\frac{a\cdot2^x+a^2-2}{2^x-1}(x\neq0)$, 其中$a$为实数. 
\begin{enumerate}[(1)]
    \item 若函数$f(x)$为奇函数, 求出$a$的值;
    \item 若函数$f(x)$在$(0,+\infty )$上是增函数, 求$a$的取值范围;
    \item 当$x\in\left[1,2\right]$时, $f(x)>1$恒成立, 求$a$的取值范围. 
    \end{enumerate}
    \begin{jieda}
    $f(x)=a+\frac{a^2+a-2}{2^x-1}$
    \begin{enumerate}
    \item $f(1)+f(-1)=0$, 可得$a=-1$或2. 再证一下. \\
    当$a=-1$时, $f(x)=-\frac{2^x+1}{2^x-1}$; 当$a=2$时, $f(x)=2\cdot\frac{2^x+1}{2^x-1}$.
    \item $g(x)=\frac{1}{2^x-1}$单调减，故$a^2+a-2\leq 0$, 即$a\in[-2,1]$，可以使得$f(x)$单调增. 
    
    \item $f(x)$必然是单调函数，因此只需其两个端点满足要求即可保证恒成立，即$f(1) > 1, f(2) > 1$, 解得$a\in(-\infty,-5)\cup(1,+\infty)$. 
    \end{enumerate}
    \end{jieda}
    \item 求函数$y=3\cdot4^x-20\cdot2^x+3\cdot4^{-x}-20\cdot2^{-x}$的最小值. 
    \begin{jieda}
    令$t=2^x+2^{-x}\in[2,+\infty)$, 则
    $y = 3t^2-20t-6 = 3(t-\frac{10}{3})^2-\frac{118}{3}$ \\
    故$y_{\min}=-\frac{118}{3}$, 当$t=\frac{10}{3}$即$2^x = \frac{1}{3}\mbox{或}3, x=log_23 \mbox{或}-log_23$时取到. 
    \end{jieda}
    \item 方程$2^x+3^x=6^x$解的个数为\blkda{1}. 
    \begin{jieda}
    $f(x)=(\frac{2}{6})^x+(\frac{3}{6})^x$严格单调, 且$f(x)\in(0,+\infty)$, 方程等价于$f(x)=1$. 
    \end{jieda}
    \item 函数$y=3^{-\abs{x-1}}-m$的图像与$x$轴有交点时, $m$的取值范围是\blkda{$(0,1]$}. 
    \begin{jieda}
    $3^{-\abs{x-1}}\in(0,1]$, 故$m\in(0,1]$. \\
    \ps 绘制$y=3^{-\abs{x-1}}$的图像的时候先考虑对称性再考虑单调性.
    \end{jieda}
    \item 当$a$为何值时, 关于$x$的方程$4^{x}-(2a+1)\cdot 2^{x}+a^2+2=0$的根一个比另一个大1? 
    \begin{jieda}
    令$t=2^x$, 化为关于$t$的方程$t^2-(2a+1)t+a^2+2=0$. \\
    $\abs{x_2-x_1}=1\iff x_2-x_1=\pm1$, 于是$\frac{t_2}{t_1}+\frac{t_1}{t_2}=2+\frac{1}{2}$. \\
    于是韦达定理结合$\Delta$解得$a=4$. 
    \end{jieda}
    
    \item 关于$x$的方程$k\cdot9^x-k\cdot3^{x+1}+6(k-5)=0$在$[0,2]$上有唯一解, 求实数$k$的取值范围. 
    \begin{jieda}
    令$t=3^x\in[1,9]$, 于是关于$t$的方程$kt^2-3kt+6(k-5)=0$在$[1,9]$上有唯一解. \\
    于是$\frac{30}{k}=t^2-3t+6$. 结合图像可得$\frac{30}{k}\in(4,60]\cup\{\frac{15}{4}\}$, 即$k\in[\frac{1}{2},\frac{15}{2})\cup\{8\}$. \\
    \ps 此处参变分离的方式值得注意！
    \end{jieda}
    

\end{enumerate}
\subsection*{对数函数、对数方程}
\begin{enumerate}
    \item 如果$\log_a2>\log_b2$, 那么$a,b$间的关系是\dotfill(\kh{D}) 
    \xx{$b>a>1$}
    {$1>b>a>0$}
    {$a>1$且$0<b<1$}
    {以上三种答案均有可能}
    \begin{jieda}
    $\iff\frac{\lg2}{\lg{a}}>\frac{\lg2}{\lg{b}}\iff\frac{1}{\lg{a}}>\frac{1}{\lg{b}}$. 
    \end{jieda}
    \item 已知函数$g(x)=f(x)-\frac{1}{f(x)}$, 其中$\log_2f(x)=2x$, $x\zr$, 则$g(x)$\dotfill(\kh{C})
    \xx{是奇函数又是减函数}{是偶函数又是增函数}{是奇函数又是增函数}{是偶函数又是减函数}
    \begin{jieda}
    $f(x)=4^x$, 则$g(x)=4^x-4^{-x}$. 
    \end{jieda}
    \item 已知$f(x)$为奇函数, 当$-1<x<0$时, $f(x)=x\lg(x+1)$, 则当$0<x<1$时, $f(x)=$\blkda[2]{$x\lg(1-x)$}. 
    \begin{jieda}
    当$x\in(0,1)$时, $-x\in(-1,0)$, 于是$f(x)=-f(-x)=-(-x)\lg(-x+1)=x\lg(-x+1)$. 
    \end{jieda}
    \item 已知函数$f(x)=\lg(1+2^x+3^x+4^x+a\cdot5^x)$对一切$x\leq 2$有意义, 求实数$a$的取值范围. 
    \begin{jieda}
    $1+2^x+3^x+4^x+a\cdot5^x>0$恒成立, 即$a>-\frac{1+2^x+3^x+4^x}{5^x}$恒成立. \\
    而$y=-\frac{1+2^x+3^x+4^x}{5^x}$在$\bR$上单调递增, 故$a>-\frac{6}{5}$($x=2$时). 
    \end{jieda}
    \item 方程$\lg{x}+x-3=0$的根所在的区间是\dotfill(\kh{B})
    \xx{$({1,2})$}{$(\frac{5}{2},\frac{11}{4})$}{$(\frac{9}{4},\frac{5}{2})$}{$(3,\frac{13}{4})$}
    \begin{jieda}
    $f(x)=\lg{x}+x-3$在$(0,+\infty)$上严格单调递增, 且$f(\frac{5}{2})f(\frac{11}{4})<0$. 
    \end{jieda}
    \item 解关于$x$的方程$\log_4(3-x)+\log_{\frac{1}{4}}(3+x)=\log_4(1-x)+\log_{\frac{1}{4}}(2x+1)$. 
    \begin{jieda}
    $\iff\log_4(3-x)-\log_4(3+x)=\log_4(1-x)-\log_4(2x+1)$. \\
    故$\frac{3-x}{3+x}=\frac{1-x}{2x+1}\land x\in(-\frac{1}{2},1)$, 解得$x=0$. 
    \end{jieda}
    \item 已知关于$x$的方程$\log_a(2x^2+x-3)-\log_a(x+4)=1+\log_a(a-1)$有实数解, 且其中之一在2与3之间. 
    \begin{enumerate}[(1)]
    \item 求实数$a$的取值范围; 
    \item 若$a$是整数, 求方程的解.
    \end{enumerate}
    \begin{jieda}
    $2x^2+x-3>0\land x+4>0\iff x\in(-4,-\frac{3}{2})\cup(1,+\infty)$. \\
    $a-1>0\land a\in(0,1)\cup(1,+\infty)\iff a>1$. 
    \begin{enumerate}[(1)]
    \item $\iff \log_a\frac{2x^2+x-3}{x+4}=\log_aa(a-1)$. \\
    故$a(a-1)=\frac{2x^2+x-3}{x+4}=2(x+4)+\frac{25}{x+4}-15$. \\
    当$x\in(2,3)$时, $x+4\in(6,7)$, 故$2(x+4)+\frac{25}{x+4}-15\in(\frac{7}{6},\frac{18}{7})$. \\
    于是$a(a-1)\in(\frac{7}{6},\frac{18}{7})\land a>1$, 故$a\in\left(\frac{3+\sqrt{51}}{6},\frac{7+\sqrt{553}}{14}\right)$. 
    
    \item $a(a-1)\in(\frac{7}{6},\frac{18}{7})\land a>1\land a\zz$, 故$a=2$. \\
    于是$2=\frac{2x^2+x-3}{x+4}$, 解得$x=\frac{1\pm\sqrt{89}}{4}$. 检验后发现符合条件. 
    \end{enumerate}
    \end{jieda}
\end{enumerate}
\section*{反函数}
\begin{enumerate}
    \item 求下列函数的反函数: \\
    \begin{enumerate*}[(1)]
    \item $y=\frac{3x-1}{2x+1}(x\neq -\frac{1}{2})$; 
    
    \item $y=\begin{cases}
    x^2-1, & 0\leq x\leq 1\\
    x^{\frac{2}{3}}, & -1\leq x<0
    \end{cases}$; 
    
    \item $y=x\abs{x}+2x$.
    \end{enumerate*}
    \begin{jieda}
    \begin{enumerate*}[(1)]
    \item $y=\frac{1+x}{3-2x}$; 
    \item $y=\begin{cases}
    \sqrt{x+1}, & -1\leq x\leq 0\\
    -x^{\frac{3}{2}}, & 0<x\leq 1
    \end{cases}$; 
    \item $y=\begin{cases}
    1-\sqrt{1-x}, & x<0\\
    \sqrt{1+x}-1, & x\geq 0
    \end{cases}$. 
    \end{enumerate*} \\
    \ps 分段函数的反函数要分段求解，第二个极易错：用$y$表示$x$的过程中，遇到开方操作时要回顾定义域确定其符号
    \end{jieda}
    \item 设$f^{-1}(x)$为$f(x) = 2^{x-2} + \frac{x}{2}, x \in [0, 2]$的反函数，则$y = f(x) + f^{-1}(x)$的最大值为\blkda{$4$}.
    \begin{jieda}
        易知$f(x)$为增函数,值域为$[\frac{1}{4}, 2]$，故$f^{-1}(x)$是$[\frac{1}{4}, 2]$上的增函数，值域为$[0,2]$。故$f(x) + f^{-1}(x)$也是$[\frac{1}{4}, 2]$上的增函数，其最大值取在端点$x = 2$处，为$2+2=4$
    \end{jieda}
    \item 已知 $f(x)=\sqrt{x-1}$, 函数 $y=f(x)$ 的反函数为 $y=f^{-1}(x)$, 若 $f^{-1}(x)-a=f(x+a)$有实数根, 则实数 $a$ 的取值范围为\blkda{$[\frac{3}{4}, +\infty)$}.
    \begin{jieda}
        $f^{-1}(x) = x^2 + 1 (x \geq 0)$. \\
        设$g(x) = f(x+a)$，则$g^{-1}(x) = f^{-1}(x)-a$，因为$g(x)$与$g^{-1}(x)$均为增函数，因此其交点必然在$y = x$上。由图象可知，$a$的取值范围形式应为$[m, +\infty)$，$m$为边界情况对应的值.边界情况为$x^2 + 1 - a = x$只有一解，此时$a = \frac{3}{4}$，故$a \in [\frac{3}{4}, +\infty)$
    \end{jieda}
    \item 已知 $f(x)=\frac{a-x}{x-a-1}$, 若函数 $y=f(x)$ 的反函数图像的对称中心是 $(-1,3)$, 则实数 $a$ 的值是\blkda{$2$}.
    \begin{jieda}
        由反函数的对称中心知$f(x)$的对称中心为$(3, -1)$，因此$a+1 = 3$，故$a = 2$
    \end{jieda}
    \item 已知函数$y=x^2-2ax+3(x\in[1,2]\cup[8,9])$有反函数, 求实数$a$的取值范围. 
    \begin{jieda}
    $a\in(-\infty,1]\cup[2,\frac{9}{2})\cup(\frac{11}{2},8]\cup[9,+\infty)$. \\ 
    \ps 根据对称轴位置分成四种情况分类讨论，使得$f(x)$满足具有反函数的要求，即一个函数值只对应一个自变量.
    \end{jieda}


\end{enumerate}